\subsubsection{Data Integration}
The Python module provides comprehensive routines for data integration, based on Jensen-Shannon Divergence. This is very useful for cross-study data.

\textbf{Key Features:}
\begin{itemize}
    \item \textbf{Calculate Jensen-Shannon Compatibility Score}: Score to evaluate how compatible two studies are.
    \item \textbf{Perform JSD Permutation Test}: Random sampling of the input data and recalculation of the JS score to test stability of the compatibility.
    \item \textbf{Family-Wise Score Contribution}: Calculates the impact of selected families on the calculated score.
\end{itemize}

\begin{enumerate}
    \item \textbf{Compute Gene Means}

        Compute per‑gene mean expression values across replicates, ignoring any \lstinline!NaN! entries.  
        The input expression matrix is converted to Fortran order for efficient passing to the underlying compiled routine.

        \textbf{Arguments:}
        \begin{itemize}
            \item \lstinline!expr (array-like)!: array of shape \lstinline!(n_reps, n_genes)! containing expression values.
        \end{itemize}

        \textbf{Returns:}
        \begin{itemize}
            \item \lstinline!np.ndarray!: array of shape \lstinline!(n_genes)! containing per‑gene mean expression values.
        \end{itemize}

        \begin{lstlisting}[language=Python, caption={Per‑Gene Mean Expression}]
means = tox_compute_gene_means(expr)
        \end{lstlisting}

    \item \textbf{Compute Residuals}

        Compute signed residuals for each gene by subtracting the per‑gene mean expression from each replicate.  
        The expression matrix is converted to Fortran order for efficient interaction with the compiled routine.

        \textbf{Arguments:}
        \begin{itemize}
            \item \lstinline!expr (array-like)!: array of shape \lstinline!(n_reps, n_genes)! containing expression values.
            \item \lstinline!means (array-like)!: array of shape \lstinline!(n_genes)! containing per‑gene mean expression values.
        \end{itemize}

        \textbf{Returns:}
        \begin{itemize}
            \item \lstinline!np.ndarray!: array of shape \lstinline!(n_reps, n_genes)! containing signed residuals.
        \end{itemize}

        \begin{lstlisting}[language=Python, caption={Signed Residual Computation}]
resid = tox_compute_residuals(expr, means)
        \end{lstlisting}


    \item \textbf{Pool Mean Expression Values}

        Pool per‑gene mean expression values across two studies and compute empirical reference points \lstinline!x_star! using the underlying compiled routine.  
        Both input vectors are flattened, concatenated internally, and NaN values are handled by the Fortran implementation.

        \textbf{Arguments:}
        \begin{itemize}
            \item \lstinline!mean_S1 (array-like)!: array of shape \lstinline!(n_genes_S1)! containing per‑gene mean expression values for study 1.
            \item \lstinline!mean_S2 (array-like)!: array of shape \lstinline!(n_genes_S2)! containing per‑gene mean expression values for study 2.
            \item \lstinline!n_points (int)!: number of reference points to compute.
        \end{itemize}

        \textbf{Returns:}
        \begin{itemize}
            \item \lstinline!dict!: \{
                \lstinline!'n_pool': int!,  
                \lstinline!'x_star': np.ndarray (n_points,)!
            \}
        \end{itemize}

        \begin{lstlisting}[language=Python, caption={Pooling Mean Expression Values}]
out = tox_pool_means(mean_S1, mean_S2, n_points)

n_pool = out["n_pool"]
x_star = out["x_star"]
        \end{lstlisting}

    \item \textbf{Pool Mean Expression Values (Expert)}

        Pool mean‑expression values using pre‑pooled and pre‑sorted arrays.  
        This expert interface assumes that pooling and sorting have already been performed externally, allowing direct computation of reference points \lstinline!x_star! without re‑allocating or re‑sorting data.

        \textbf{Arguments:}
        \begin{itemize}
            \item \lstinline!pooled_means (array-like)!: 1D array of shape \lstinline!(pool_size)!, containing pre‑pooled mean expression values.
            \item \lstinline!pooled_perm (array-like)!: permutation vector of shape \lstinline!(pool_size)! that sorts \lstinline!pooled_means!.
            \item \lstinline!n_points (int)!: number of reference points to compute.
        \end{itemize}

        \textbf{Returns:}
        \begin{itemize}
            \item \lstinline!dict!: \{
                \lstinline!'n_pool': int!,  
                \lstinline!'x_star': np.ndarray (n_points,)!
            \}
        \end{itemize}

        \begin{lstlisting}[language=Python, caption={Pooling Means (Expert Interface)}]
out = tox_pool_means_expert(
    pooled_means,
    pooled_perm,
    n_points
)

n_pool = out["n_pool"]
x_star = out["x_star"]
        \end{lstlisting}

    \item \textbf{Compute Neighborhood Size}

        Compute the number of neighbors to use for neighborhood‑based residual analysis.  
        The calculation mirrors the Fortran routine: it enforces minimum and maximum limits, accounts for the number of reference points, and excludes genes with \lstinline!NaN! mean expression values.

        \textbf{Arguments:}
        \begin{itemize}
            \item \lstinline!n_pool (int)!: total number of pooled mean‑expression values across both studies.
            \item \lstinline!n_points (int)!: number of reference points.
            \item \lstinline!n_genes_S (int)!: number of genes in the current study.
            \item \lstinline!mean_S (array-like)!: array of shape \lstinline!(n_genes_S)! containing per‑gene mean expression values.
            \item \lstinline!desired_size (int)!: optional desired neighborhood size (default \lstinline!0!, meaning “use Fortran default logic”).
        \end{itemize}

        \textbf{Returns:}
        \begin{itemize}
            \item \lstinline!int!: computed neighborhood size.
        \end{itemize}

        \begin{lstlisting}[language=Python, caption={Neighborhood Size Computation}]
n_neighbors = tox_calc_neighborhood_size(
    n_pool,
    n_points,
    n_genes_S,
    mean_S,
    desired_size=0
)
        \end{lstlisting}

    \item \textbf{Construct Neighborhoods}

        Construct neighborhood‑based residual sets by selecting, for each reference point, the closest genes in mean‑expression space.  
        The neighborhood size is computed automatically unless a desired size is provided.  
        Residuals and indices are returned in arrays sized to the actual number of neighbors.

        \textbf{Arguments:}
        \begin{itemize}
            \item \lstinline!x_star (array-like)!: array of shape \lstinline!(n_points)! containing reference mean‑expression values.
            \item \lstinline!mean_S (array-like)!: array of shape \lstinline!(n_genes_S)! with per‑gene mean expression values.
            \item \lstinline!resid_S (array-like)!: array of shape \lstinline!(n_reps_S, n_genes_S)! containing signed residuals.
            \item \lstinline!n_pool (int)!: total number of pooled mean‑expression values across studies.
            \item \lstinline!desired_n_neighbors (int)!: optional desired neighborhood size (default \lstinline!0!, meaning “compute automatically”).
        \end{itemize}

        \textbf{Returns:}
        \begin{itemize}
            \item \lstinline!dict!: \{
                \lstinline!'neighborhood_residuals': np.ndarray (n_reps_S, actual_n_neighbors, n_points)!,  
                \lstinline!'neighborhood_indices': np.ndarray (actual_n_neighbors, n_points)!
            \}
        \end{itemize}

        \begin{lstlisting}[language=Python, caption={Neighborhood Construction}]
out = tox_construct_neighborhoods(
    x_star,
    mean_S,
    resid_S,
    n_pool,
    desired_n_neighbors=0
)

neigh_res = out["neighborhood_residuals"]
neigh_idx = out["neighborhood_indices"]
        \end{lstlisting}

    \item \textbf{Determine Shared Residual Range (Expert)}

        Compute the shared residual range \(R\) from a pre‑pooled and pre‑sorted residual vector using a percentile threshold.

        \textbf{Arguments:}
        \begin{itemize}
            \item \lstinline!residual_pool (array-like)!: 1D array of shape \lstinline!(pool_size)!, typically \lstinline!(n_reps_S1 + n_reps_S2) * n_neighbors * n_points!.
            \item \lstinline!residual_pool_perm (array-like)!: permutation vector of shape \lstinline!(pool_size)! that sorts \lstinline!residual_pool!.
            \item \lstinline!residual_range_quantile (float)!: percentile used to determine the residual range (default \lstinline!95.0!).
        \end{itemize}

        \textbf{Returns:}
        \begin{itemize}
            \item \lstinline!float!: shared residual range \(R\)
        \end{itemize}

        \begin{lstlisting}[language=Python, caption={Shared Residual Range (Expert)}]
R = tox_determine_shared_residual_range_expert(
    residual_pool,
    residual_pool_perm,
    residual_range_quantile=95.0
)
        \end{lstlisting}
 
    \item \textbf{Determine Shared Residual Range}

        Compute the shared residual range \(R\) from two residual matrices.  
        Both residual tensors are flattened (absolute values), pooled, and the percentile threshold is applied to determine the symmetric residual range \([-R, R]\).

        \textbf{Arguments:}
        \begin{itemize}
            \item \lstinline!neighborhood_residuals_S1 (array-like)!: array of shape \lstinline!(n_reps_S1, n_neighbors, n_points)!.
            \item \lstinline!neighborhood_residuals_S2 (array-like)!: array of shape \lstinline!(n_reps_S2, n_neighbors, n_points)!.
            \item \lstinline!residual_range_quantile (float)!: percentile used to determine the residual range (default \lstinline!95.0!).
        \end{itemize}

        \textbf{Returns:}
        \begin{itemize}
            \item \lstinline!float!: shared residual range \(R\)
        \end{itemize}

        \begin{lstlisting}[language=Python, caption={Shared Residual Range}]
R = tox_determine_shared_residual_range(
    neighborhood_residuals_S1,
    neighborhood_residuals_S2,
    residual_range_quantile=95.0
)
        \end{lstlisting}

    \item \textbf{Build Residual Histograms}

        Build histogram counts and probability mass functions (PMFs) for one study, using a shared residual range and a fixed number of bins.  
        Residuals are grouped per reference point, NaNs are ignored, and both absolute counts and normalized PMFs are returned.

        \textbf{Arguments:}
        \begin{itemize}
            \item \lstinline!neighborhood_residuals (array-like)!: array of shape \lstinline!(n_reps, n_neighbors, n_points)! containing residuals.
            \item \lstinline!shared_residual_range (float)!: symmetric residual range \([-R, R]\) used for histogram binning.
            \item \lstinline!n_bins (int)!: number of equally sized histogram bins.
        \end{itemize}

        \textbf{Returns:}
        \begin{itemize}
            \item \lstinline!dict!: \{
                \lstinline!'counts': np.ndarray (n_points, n_bins)!,  
                \lstinline!'pmf': np.ndarray (n_points, n_bins)!,  
                \lstinline!'included_n_residuals': np.ndarray (n_points,)!
            \}
        \end{itemize}

        \begin{lstlisting}[language=Python, caption={Residual Histogram Construction}]
out = tox_build_residual_histograms(
    neighborhood_residuals,
    shared_residual_range,
    n_bins
)

counts = out["counts"]
pmf = out["pmf"]
included = out["included_n_residuals"]
        \end{lstlisting}

    \item \textbf{Build Residual Histograms (Filtered)}

        Build histogram counts and probability mass functions (PMFs) for one study, using a shared residual range and a fixed number of bins.  
        A boolean neighbor mask selects which neighbors contribute to each reference point’s histogram.

        \textbf{Arguments:}
        \begin{itemize}
            \item \lstinline!neighborhood_residuals (array-like)!: array of shape \lstinline!(n_reps, n_neighbors, n_points)! containing residuals.
            \item \lstinline!shared_residual_range (float)!: symmetric residual range \([-R, R]\) used for histogram binning.
            \item \lstinline!n_bins (int)!: number of histogram bins.
            \item \lstinline!neighbor_mask (array-like)!: boolean mask of shape \lstinline!(n_neighbors, n_points)! selecting neighbors to include.
        \end{itemize}

        \textbf{Returns:}
        \begin{itemize}
            \item \lstinline!dict!: \{
                \lstinline!'counts': np.ndarray (n_points, n_bins)!,  
                \lstinline!'pmf': np.ndarray (n_points, n_bins)!,  
                \lstinline!'included_n_residuals': np.ndarray (n_points,)!
            \}
        \end{itemize}

        \begin{lstlisting}[language=Python, caption={Filtered Residual Histogram Construction}]
out = tox_build_residual_histograms_filtered(
    neighborhood_residuals,
    shared_residual_range,
    n_bins,
    neighbor_mask
)

counts = out["counts"]
pmf = out["pmf"]
included = out["included_n_residuals"]
        \end{lstlisting}

    \item \textbf{Compute Divergence per Reference Point}

        Compute the Jensen–Shannon divergence for each reference point using two probability mass function (PMF) matrices.  
        Each row corresponds to a reference point, and each column to a histogram bin.

        \textbf{Arguments:}
        \begin{itemize}
            \item \lstinline!pmf_S1 (array-like)!: array of shape \lstinline!(n_points, n_bins)! containing PMFs for study 1.
            \item \lstinline!pmf_S2 (array-like)!: array of shape \lstinline!(n_points, n_bins)! containing PMFs for study 2.
        \end{itemize}

        \textbf{Returns:}
        \begin{itemize}
            \item \lstinline!np.ndarray!: array of shape \lstinline!(n_points,)! containing per‑point Jensen–Shannon divergences.
        \end{itemize}

        \begin{lstlisting}[language=Python, caption={Per‑Point JSD Computation}]
jsd = tox_compute_divergence_per_reference_point(
    pmf_S1,
    pmf_S2
)
        \end{lstlisting}


    \item \textbf{Compute Weighted Global Divergence}

        Compute the weighted global Jensen–Shannon divergence using per‑point JSD values and the number of included residuals from both studies.  
        Weights are proportional to the total number of included residuals per reference point.

        \textbf{Arguments:}
        \begin{itemize}
            \item \lstinline!js_divergences (array-like)!: array of shape \lstinline!(n_points)! containing per‑point JSD values.
            \item \lstinline!included_S1 (array-like)!: array of shape \lstinline!(n_points)! with included residual counts for study 1.
            \item \lstinline!included_S2 (array-like)!: array of shape \lstinline!(n_points)! with included residual counts for study 2.
        \end{itemize}

        \textbf{Returns:}
        \begin{itemize}
            \item \lstinline!dict!: \{
                \lstinline!'global_js_divergence': float!,  
                \lstinline!'weights': np.ndarray (n_points,)!
            \}
        \end{itemize}

        \begin{lstlisting}[language=Python, caption={Weighted Global JSD}]
out = tox_compute_weighted_global_divergence(
    js_divergences,
    included_S1,
    included_S2
)

global_jsd = out["global_js_divergence"]
weights = out["weights"]
        \end{lstlisting}

    \item \textbf{Permutation Test for Global JSD}

        Estimate how likely the observed global Jensen–Shannon divergence is to occur by chance under the null hypothesis that both studies are exchangeable.  
        Residuals from both studies are repeatedly pooled and shuffled, histograms are rebuilt, and global JSD values are recomputed for each permutation.

        \textbf{Arguments:}
        \begin{itemize}
            \item \lstinline!neighborhood_residuals_S1 (array-like)!: array of shape \lstinline!(n_reps, n_neighbors, n_points)! for study 1.
            \item \lstinline!neighborhood_residuals_S2 (array-like)!: array of shape \lstinline!(n_reps, n_neighbors, n_points)! for study 2.
            \item \lstinline!global_jsd_observed (float)!: observed global Jensen–Shannon divergence.
            \item \lstinline!n_bins (int)!: number of histogram bins.
            \item \lstinline!shared_residual_range (float)!: shared residual range \([-R, R]\) used for histogram binning.
            \item \lstinline!n_permutations (int)!: number of permutations to perform.
            \item \lstinline!random_seed (int)!: seed for reproducible shuffling (default \lstinline!42!).
        \end{itemize}

        \textbf{Returns:}
        \begin{itemize}
            \item \lstinline!dict!: \{
                \lstinline!'jsd_null': np.ndarray (n_permutations,)!,  
                \lstinline!'p_value': float!
            \}
        \end{itemize}

        \begin{lstlisting}[language=Python, caption={Permutation Test for Global JSD}]
out = gjct_permutation_test(
    neighborhood_residuals_S1,
    neighborhood_residuals_S2,
    global_jsd_observed,
    n_bins,
    shared_residual_range,
    n_permutations,
    random_seed=42
)

jsd_null = out["jsd_null"]
p_value = out["p_value"]
        \end{lstlisting}

    \item \textbf{Permutation Test for Global JSD (Filtered)}

        Estimate how likely the observed global Jensen–Shannon divergence is to occur by chance under the null hypothesis that both studies are exchangeable.  
        Neighbor masks allow excluding specific neighbors for each reference point in each study during histogram construction.

        \textbf{Arguments:}
        \begin{itemize}
            \item \lstinline!neighborhood_residuals_S1 (array-like)!: array of shape \lstinline!(n_reps, n_neighbors, n_points)! for study 1.
            \item \lstinline!neighborhood_residuals_S2 (array-like)!: array of shape \lstinline!(n_reps, n_neighbors, n_points)! for study 2.
            \item \lstinline!global_jsd_observed (float)!: observed global Jensen–Shannon divergence.
            \item \lstinline!n_bins (int)!: number of histogram bins.
            \item \lstinline!shared_residual_range (float)!: shared residual range \([-R, R]\) used for histogram binning.
            \item \lstinline!n_permutations (int)!: number of permutations to perform.
            \item \lstinline!neighbor_mask_S1 (array-like)!: boolean mask \lstinline!(n_neighbors, n_points)! selecting neighbors for study 1.
            \item \lstinline!neighbor_mask_S2 (array-like)!: boolean mask \lstinline!(n_neighbors, n_points)! selecting neighbors for study 2.
            \item \lstinline!random_seed (int)!: seed for reproducible shuffling (default \lstinline!42!).
        \end{itemize}

        \textbf{Returns:}
        \begin{itemize}
            \item \lstinline!dict!: \{
                \lstinline!'jsd_null': np.ndarray (n_permutations,)!,  
                \lstinline!'p_value': float!
            \}
        \end{itemize}

        \begin{lstlisting}[language=Python, caption={Filtered Permutation Test for Global JSD}]
out = gjct_permutation_test_filtered(
    neighborhood_residuals_S1,
    neighborhood_residuals_S2,
    global_jsd_observed,
    n_bins,
    shared_residual_range,
    n_permutations,
    neighbor_mask_S1,
    neighbor_mask_S2,
    random_seed=42
)

jsd_null = out["jsd_null"]
p_value = out["p_value"]
        \end{lstlisting}

    \item \textbf{Compute Family-Level JSD}

        Compute the family-level Jensen–Shannon divergence for a given paralog family.  
        This routine builds family-specific neighbor masks, constructs masked residual histograms, computes per‑point JSD values, and aggregates them into a weighted global divergence.

        \textbf{Arguments:}
        \begin{itemize}
            \item \lstinline!family_idx (int)!: index of the gene family to analyze.
            \item \lstinline!gene_to_family_S1 (array-like)!: \lstinline!(n_genes_S1)! int32 vector mapping each gene in study 1 to a family.
            \item \lstinline!gene_to_family_S2 (array-like)!: \lstinline!(n_genes_S2)! int32 vector mapping each gene in study 2 to a family.
            \item \lstinline!neighborhood_residuals_S1 (array-like)!: array of shape \lstinline!(n_reps_S1, n_neighbors, n_points)!.
            \item \lstinline!neighborhood_residuals_S2 (array-like)!: array of shape \lstinline!(n_reps_S2, n_neighbors, n_points)!.
            \item \lstinline!neighborhood_genes_S1 (array-like)!: array of shape \lstinline!(n_neighbors, n_points)! with gene indices for study 1.
            \item \lstinline!neighborhood_genes_S2 (array-like)!: array of shape \lstinline!(n_neighbors, n_points)! with gene indices for study 2.
            \item \lstinline!n_bins (int)!: number of histogram bins.
            \item \lstinline!shared_residual_range (float)!: shared residual range \([-R, R]\) used for histogram construction.
        \end{itemize}

        \textbf{Returns:}
        \begin{itemize}
            \item \lstinline!dict!: \{
                \lstinline!'js_divergences': np.ndarray (n_points,)!,\\
                \lstinline!'included_n_reps_S1': np.ndarray (n_points,)!,\\
                \lstinline!'included_n_reps_S2': np.ndarray (n_points,)!,\\
                \lstinline!'total_included_n_reps': int!,\\
                \lstinline!'global_js_divergence': float!,\\
                \lstinline!'weights': np.ndarray (n_points,)!,\\
                \lstinline!'ierr': int!
            \}
        \end{itemize}

        \begin{lstlisting}[language=Python, caption={Family-Level JSD Computation}]
out = fjct_compute_jsd(
    family_idx,
    gene_to_family_S1,
    gene_to_family_S2,
    neighborhood_residuals_S1,
    neighborhood_residuals_S2,
    neighborhood_genes_S1,
    neighborhood_genes_S2,
    n_bins,
    shared_residual_range
)

jsd = out["js_divergences"]
global_jsd = out["global_js_divergence"]
weights = out["weights"]
        \end{lstlisting}
\end{enumerate}